\section{The Art of Arts: \work{Regula Pastoralis} and the Theology of Office}

\subsection{Occasion and design}

John, Archbishop of Ravenna, wrote kindly to the new pope in the autumn of 590
to reprove him for having tried to hide from the pastoral burden. Gregory
answered not with a brief apology, but with a treatise that became the classic
handbook of Western Christian leadership for a thousand years: the \work{Liber
Regulae Pastoralis}, or \work{Book of Pastoral Rule}. Composed at the very
beginning of his episcopate (as he later told Leander in \work{Ep.}\ V.49) and
drawing on an outline he had already sketched in \work{Moralia in Iob}
XXX.12--13, the book set out to prove why any sane man trembles before the
governance of souls.

Gregory openly acknowledged his debt to Saint Gregory of Nazianzus's second
oration (\work{Regula Pastoralis} III, Prologue), which had likewise defended a
reluctant flight from priesthood. Yet Pope Gregory organized his work with
distinctive Roman clarity into four ascending steps: how a man should come to
pastoral office (Part I); how he should live once set in it (Part II); how he
should teach and admonish the varied souls entrusted to him (Part III); and
how, after doing all things well, he must return daily to his own weakness lest
pride ruin what virtue built (Part IV).

\subsection{The burden and the balance of the shepherd}

The opening chapter states the governing conviction of the entire book:

\begin{quote}
No one presumes to teach an art till he has first, with intent meditation,
learnt it. What rashness is it, then, for the unskilful to assume pastoral
authority, since the government of souls is the art of arts! For who can be
ignorant that the sores of the thoughts of men are more occult than the sores
of the bowels? And yet how often do men who have no knowledge whatever of
spiritual precepts fearlessly profess themselves physicians of the heart,
though those who are ignorant of the effect of drugs blush to appear as
physicians of the flesh! (\work{Regula Pastoralis} I.1)
\end{quote}

Part I cuts cleanly in two directions. First, Gregory warns off the ambitious,
the unlearned, and the hypocrites who study Scripture eagerly yet trample its
precepts in their conduct---shepherds who drink clear water and foul the rest
with their feet so that the flock drinks mud (I.2, citing Ezek.\ 34:18--19).
Christ himself fled when the crowd tried to make him king and went willingly
only to the cross, teaching his followers to fear prosperity and bear adversity
(I.3). Office scatters the mind across a hundred outward concerns until a ruler
``knows how to think of many things, while itself it knows not'' (I.4). Yet
Gregory refuses to let monastic escapism win the argument. Those who truly
possess the virtues needed to guide others sin if they refuse the call out of
love for their own ease: ``If, then, the care of feeding is the proof of
loving, whosoever abounds in virtues, and yet refuses to feed the flock of God,
is convicted of not loving the chief Shepherd'' (I.5, citing John 21:17).

Part II sketches the interior architecture of the bishop who has accepted that
paradox. Gregory's ideal pastor is never a mere administrator and never a
detached mystic; he must hold together action and contemplation in daily
tension. The pastor must be ``pure in thought, in action chief; discreet in
keeping silence, profitable in speech; a near neighbour to every one in
sympathy, exalted above all in contemplation; a familiar friend of good livers
through humility, unbending against the vices of evil-doers through zeal for
righteousness'' (\work{Regula Pastoralis} II.1; echoed verbatim in \work{Ep.}\
I.25).

Two chapters in Part II reveal Gregory at the summit of his spiritual writing.
In II.4, meditating on the golden bells and pomegranates sewn into the hem of
Aaron's priestly vestment (Exod.\ 28:33--35), he explains that the bells are
the sound of preaching---for the priest who enters the sanctuary silent dies
before God---while the pomegranates are the unity of faith, many seeds bound
within one rind: ``Join ye pomegranates to bells, that in all ye say ye may
with cautious watchfulness keep the unity of the faith.'' In II.5, he insists
that true contemplation does not alienate the pastor from human messiness; it
deepens his compassion. Saint Paul was caught up to the third heaven, yet bent
his gaze down to counsel married couples on the bed of their weakness; Moses
went continually into the tabernacle to speak with God and came back out to
bear the burdens of a quarrelsome people; Christ spent the night in prayer on
the mountain and healed the sick in the cities. ``Charity then rises
wonderfully to high things when it is compassionately drawn to the low things
of neighbours.'' Even more striking is Gregory's insistence in II.6 on the
radical equality of human nature: ``it is clear that nature produced all men
equal; but, through variation in the order of their merits, guilt puts some
below others.'' God gave Adam and Noah dominion over beasts, not over fellow
men; therefore ``to wish to be feared by one's equal is to be proud against
nature,'' and ``supreme rule, then, is ordered well, when he who presides lords
it over vices, rather than over his brethren'' (\work{Regula Pastoralis} II.6).

\subsection{Pastoral psychology and the snare of virtue}

Part III---by far the longest section of the book---displays Gregory's
unrivaled eye for the human heart. Because no single medicine heals every
disease and ``the bread which invigorates the life of the strong kills little
children,'' the preacher must handle his hearers like the strings of a harp:
``the strings render back a consonant modulation, that they are struck indeed
with one quill, but not with one kind of stroke'' (\work{Regula Pastoralis}
III, Prologue). Gregory then works systematically through thirty-six paired
conditions of soul (Chapter 1 and the admonitions that follow): men and women,
young and old, poor and rich, joyful and sad, servants and masters, the
impudent and the bashful, the forward and the fainthearted, the obstinate and
the fickle, the married and the single, those who bewail sins without forsaking
them and those who forsake them without bewailing them, those who do evil
secretly and good publicly and those who hide their good deeds while allowing
evil to be thought of them. Each pair receives a tailored spiritual diagnosis
grounded in Scripture and lived observation.

Part IV brings the pastor back from the pulpit to his own conscience in a
single piercing chapter. The deadliest danger awaits not the negligent ruler,
but the successful one: ``before the eyes of the just judge the memory of
virtue is a pitfall of the soul'' (\work{Regula Pastoralis} IV). God often
leaves even holy pastors struggling with small lingering faults so that they
cannot boast of their great achievements. Gregory ends by turning the mirror
directly on himself:

\begin{quote}
See now, good man, how, compelled by the necessity laid upon me by thy reproof,
being intent on shewing what a Pastor ought to be, I have been as an
ill-favoured painter pourtraying a handsome man; and how I direct others to the
shore of perfection, while myself still tossed among the waves of
transgressions. But in the shipwreck of this present life sustain me, I beseech
thee, by the plank of thy prayer, that, since my own weight sinks me down, the
hand of thy merit may raise me up. (\work{Regula Pastoralis} IV)
\end{quote}

\subsection{Immediate and enduring reception}

The treatise conquered Christendom within Gregory's own lifetime. Bishop
Licinianus of Cartagena in Spain wrote to Gregory praising the book as ``a
palace of all virtues'' attested by Hilary, Ambrose, Augustine, and Gregory of
Nazianzus, while adding with refreshing candor that if Gregory's standard were
enforced strictly, scarcely anyone would dare to be ordained (\work{Ep.}\
II.54). Emperor Maurice obtained a copy and had Patriarch Anastasius of
Antioch translate it into Greek (reported by Barmby from \work{Ep.}\ XII.24).
Two centuries later, Charlemagne's reforming councils of 813 (Mainz, Rheims,
Tours, and Ch\^alon-sur-Sa\^one) and the Council of Aachen (836) required
bishops to study it beside Scripture and the canons; Archbishop Hincmar of
Rheims reported that bishops received it in their hands at the altar during
consecration; and in late ninth-century England, King Alfred the Great
translated it into Old English so that every bishopric in his kingdom might
keep a copy (Barmby, Preface to \work{Regula Pastoralis}).
